\section{Introduction}

Modern compilers do most of their interesting work in the middle end, on a
typed, static single assignment intermediate representation that sits between the
language frontend and the machine specific backend. LLVM is the dominant
production infrastructure for this style of work \cite{lattner2004llvm}, and its
new pass manager is the interface through which analyses and transforms are
written today \cite{llvm_writing_pass}.

This project builds a small but complete slice of that world from the outside. I
wrote a plugin that a stock \texttt{opt} binary loads, containing analysis passes
that measure the IR and transformation passes that rewrite it, wired into a
configurable pipeline. The goal was never to beat LLVM's own optimizer. It was to
demonstrate command of the infrastructure: how analyses cache and invalidate,
how transforms report what they preserve, how a pipeline reaches a fixed point,
and how to prove that a rewrite is correct rather than merely plausible.

The guiding rule throughout was honesty. Every metric in this report comes from a
recorded run on the target machine; there are no hand typed numbers in the
document source. Every transform is gated by a differential execution check that
runs the compiled kernel before and after and compares its output, so a rewrite
that changes observable behavior fails the build regardless of how much IR it
removed. Where the framework does little, the report says so plainly, including a
kernel on which it correctly does nothing at all.
